% META: {"id":"1SPE-DERIVATION-GLOBAL-FR-R2","chapitre":"1SPE-DERIVATION-GLOBAL","type_objet":"remediation","capacites_codes":["R2"],"status":"generated"}

%---------------------------------------------------------------
% FICHE DE REMISE À NIVEAU — R2 : Équation de la tangente
%---------------------------------------------------------------

\subsection*{Rappel — Équation de la tangente}

La tangente à la courbe de $f$ au point d'abscisse $a$ a pour équation :
\[
  y = f'(a)(x - a) + f(a).
\]
Elle passe par le point $(a\,;\,f(a))$ et a pour coefficient directeur $f'(a)$.

%---------------------------------------------------------------

\begin{exercice}{1SPE-DERGLOBAL-FR-R2-EX1}{1}{6}
On sait que $f(2) = 3$ et $f'(2) = 5$.
\begin{enumerate}
  \item Écrire l'équation de la tangente $T$ à la courbe de $f$ au point d'abscisse $2$.
  \item Calculer $T(2)$ et vérifier que le point $(2\,;\,3)$ est bien sur $T$.
  \item Calculer $T(3)$.
\end{enumerate}
\end{exercice}

% BEGIN-VERIFY
% from sympy import symbols, expand, simplify
% x = symbols('x')
% T = expand(5*(x-2) + 3)
% assert T == 5*x - 7
% assert T.subs(x, 2) == 3
% assert T.subs(x, 3) == 8
% END-VERIFY

\begin{corrige}{1SPE-DERGLOBAL-FR-R2-EX1}

\commentaireMarge{On applique la formule $y = f'(a)(x-a)+f(a)$ avec $a=2$, $f(2)=3$, $f'(2)=5$.}

\textbf{Question 1.}
$T : y = 5(x - 2) + 3 = 5x - 10 + 3 = 5x - 7.$

\textbf{Question 2.}
$T(2) = 5 \times 2 - 7 = 10 - 7 = 3 = f(2).$ \checkmark Le point $(2\,;\,3)$ est bien sur $T$.

\textbf{Question 3.}
$T(3) = 5 \times 3 - 7 = 15 - 7 = 8.$

\end{corrige}

%---------------------------------------------------------------

\begin{exercice}{1SPE-DERGLOBAL-FR-R2-EX2}{1}{6}
Soit $f(x) = x^2$. On admet que $f'(x) = 2x$.
\begin{enumerate}
  \item Calculer $f(-1)$ et $f'(-1)$.
  \item Écrire l'équation de la tangente à la courbe de $f$ au point d'abscisse $-1$.
  \item Vérifier que la tangente coupe l'axe des ordonnées en $(0\,;\,-1)$.
\end{enumerate}
\end{exercice}

% BEGIN-VERIFY
% from sympy import symbols, expand, simplify
% x = symbols('x')
% f = x**2
% fp = 2*x
% assert f.subs(x, -1) == 1
% assert fp.subs(x, -1) == -2
% T = expand(fp.subs(x,-1)*(x-(-1)) + f.subs(x,-1))
% assert simplify(T - (-2*x - 1)) == 0
% assert T.subs(x, 0) == -1
% END-VERIFY

\begin{corrige}{1SPE-DERGLOBAL-FR-R2-EX2}

\textbf{Question 1.}
$f(-1) = (-1)^2 = 1$ et $f'(-1) = 2 \times (-1) = -2$.

\textbf{Question 2.}
$T : y = f'(-1)(x - (-1)) + f(-1) = -2(x + 1) + 1 = -2x - 2 + 1 = -2x - 1.$

\textbf{Question 3.}
$T(0) = -2 \times 0 - 1 = -1.$ La tangente coupe bien l'axe des ordonnées en $(0\,;\,-1)$. \checkmark

\end{corrige}

%---------------------------------------------------------------

\begin{exercice}{1SPE-DERGLOBAL-FR-R2-EX3}{1}{7}
On donne la courbe $\mathcal{C}_f$. On lit que $f(1) = 4$ et que la tangente en $(1\,;\,4)$ passe par le point $(3\,;\,0)$.
\begin{enumerate}
  \item Calculer la pente de la tangente.
  \item En déduire $f'(1)$.
  \item Écrire l'équation de la tangente.
\end{enumerate}
\end{exercice}

% BEGIN-VERIFY
% from sympy import Rational, symbols, expand, simplify
% x = symbols('x')
% # pente = (0-4)/(3-1) = -2
% pente = Rational(0-4, 3-1)
% assert pente == -2
% T = expand(-2*(x-1) + 4)
% assert simplify(T - (-2*x + 6)) == 0
% assert T.subs(x, 1) == 4
% assert T.subs(x, 3) == 0
% END-VERIFY

\begin{corrige}{1SPE-DERGLOBAL-FR-R2-EX3}

\commentaireMarge{La pente d'une droite passant par $(x_1;y_1)$ et $(x_2;y_2)$ est $\frac{y_2-y_1}{x_2-x_1}$.}

\textbf{Question 1.}
$\text{pente} = \dfrac{0 - 4}{3 - 1} = \dfrac{-4}{2} = -2.$

\textbf{Question 2.}
$f'(1) = -2$ (la pente de la tangente est le nombre dérivé).

\textbf{Question 3.}
$T : y = -2(x - 1) + 4 = -2x + 2 + 4 = -2x + 6.$

Vérification : $T(3) = -6 + 6 = 0$ \checkmark et $T(1) = -2 + 6 = 4 = f(1)$ \checkmark.

\end{corrige}
